\section{Full-data tuning and gradient-based metrics}\label{sec:retune}

We compare hyperparameter selection on a 6{,}000-row subsample with selection on all
training rows, and coordinate search with a metric read off the gradients of a cubic surrogate
(Table~\ref{tab:retune}).

\begin{table}[t]
\centering\small
\resizebox{\textwidth}{!}{\begin{tabular}{lcccccc}
\toprule
 & \multicolumn{4}{c}{component rel.\ $L^2$ [\%]} & radiance & reflectance \\
\cmidrule(lr){2-5}
kernel & $Y_1$ & $Y_2$ & $Y_3$ & $Y_4$ & rel.\ $L^2$ [\%] & med.\ $|\hat\rho-\rho|$ [\%] \\
\midrule
isotropic, tuned on 6{,}000 rows & 0.08$\pm$0.01 & 0.42$\pm$0.03 & 0.20$\pm$0.00 & 2.62$\pm$0.18 & 0.239$\pm$0.016 & 0.141$\pm$0.018 \\
isotropic, tuned on every row & 0.08$\pm$0.00 & 0.36$\pm$0.02 & 0.17$\pm$0.01 & 2.58$\pm$0.14 & 0.210$\pm$0.008 & 0.107$\pm$0.006 \\
per-input search, tuned on 6{,}000 rows & 0.08$\pm$0.01 & 0.16$\pm$0.02 & 0.08$\pm$0.01 & 2.18$\pm$0.16 & 0.095$\pm$0.009 & 0.036$\pm$0.004 \\
per-input search, re-tuned on every row & 0.08$\pm$0.00 & 0.16$\pm$0.02 & 0.08$\pm$0.01 & 2.09$\pm$0.12 & 0.096$\pm$0.009 & 0.036$\pm$0.004 \\
per-input from the cubic's gradients, $p=1$ & 0.34$\pm$0.01 & 0.19$\pm$0.02 & 0.14$\pm$0.01 & 2.17$\pm$0.14 & 0.159$\pm$0.007 & 0.078$\pm$0.004 \\
per-input from the cubic's gradients, $p=1/2$ & 0.11$\pm$0.00 & 0.19$\pm$0.02 & 0.09$\pm$0.01 & 2.27$\pm$0.15 & 0.125$\pm$0.008 & 0.055$\pm$0.004 \\
per-input from the cubic's gradients, $p=1/4$ & 0.08$\pm$0.00 & 0.25$\pm$0.02 & 0.14$\pm$0.01 & 2.38$\pm$0.20 & 0.157$\pm$0.015 & 0.077$\pm$0.013 \\
\bottomrule
\end{tabular}
}
\caption{Where the kernel's hyperparameters are chosen, and how its metric is obtained, at the
full training block and the same ten splits. The first pair tunes smoothness, scale and nugget
on a $6{,}000$-row subsample and on every training row; the last three rows replace the coordinate
search by a metric read off the cubic surrogate's gradients, at three exponents. Selecting the
exponent on validation, per split, chooses $p=1/2$ on all ten, so that row would repeat the
$p=1/2$ row and is omitted. Mean $\pm$ standard deviation over the ten splits.}
\label{tab:retune}
\end{table}

\paragraph{Tuning with every row in the solve.} The kernel families of
Section~\ref{sec:methods} choose smoothness, scale and nugget by the validation error of a solve
on a $6{,}000$-row subsample and refit the winner on all $18{,}884$ rows. Tuning with the full
solve instead lowers the isotropic kernel's radiance error from $0.239\%$ to $0.210\%$ and its
median reflectance error from $0.141$ to $0.107$ points, better on nine of ten splits, with the
ratio averaging $0.882$ and never above $1.001$. The selected scales differ: over the forty
component-split cells the full block picks a shorter length scale than the subsample in $26$,
the same in $12$ and a longer one in $2$. A smaller tuning fit prefers more smoothing.

Below $6{,}000$ rows the tuning subsample is the whole training block, so those rungs are
already fully tuned; at $8{,}000$ rows it is three quarters of the block and only at the top rung
is it a third. Applying the measured ratio to the top rung of Table~\ref{tab:curves} takes the
kernel's last step on radiance from $-0.183$ to $-0.334$, against a slope of $-0.305$ over the
five lower rungs, and on the two flux components from $-0.187$ and $-0.171$ to $-0.352$ and
$-0.358$, against $-0.437$ and $-0.468$. Most of the bend of Section~\ref{sec:scaling} is the
tuning protocol, and with full tuning the kernel's margin over the network at the full block is
$1.8$ rather than $1.6$.

Full-data tuning barely changes the input-scaled kernel: $0.096\%$ against $0.095\%$, better on
four of ten splits, ratio $1.005$. Its coordinate search already arrives at the scale the full
block selects.

\paragraph{Gradient-based metric.} The coordinate search costs a sweep over the six inputs at
five multipliers, twice, per component and split. The alternative fixes the metric before any
kernel is fitted: the multiplier of coordinate $j$ is the root-mean-square gradient of the fitted
cubic surrogate along $j$ over the training rows, raised to a power $p$, normalized to geometric
mean one and clipped to $[1/16, 4]$, the range of the multipliers the search selected, so that
weakly varying coordinates receive longer length scales. It costs one cubic fit and six finite differences.

At $p=1$ the radiance error is $0.159\%$, below the isotropic kernel on every split, by a factor
$0.669$ on average, and above the searched metric on every split, by a factor $1.680$. The
dependence on $p$ is not monotone: $p=1/2$ gives $0.125\%$ and $p=1/4$ gives $0.157\%$. Choosing
$p$ like every other hyperparameter here, by the lowest mean validation error over the four
components, selects $p=1/2$ on all ten splits. At that exponent the gradient metric closes $79\%$
of the gap between the isotropic and searched kernels on radiance and $82\%$ on the median
reflectance error, for one cubic fit and a choice among three numbers.

The components prefer different exponents: $Y_2$ and $Y_3$ are best at $1/2$, the albedo at $1$
($2.17\%$ against $2.27\%$), and the path radiance at $1/4$. On path radiance the isotropic kernel
already reads $0.084\%$ and the searched metric $0.085\%$, while the proportional metric reads
$0.34\%$, four times the isotropic error. The gradient of the surrogate along a coordinate
measures how fast the map moves with it, not how much that coordinate helps predict, and on this
component the two disagree. Softening the exponent to $1/4$ brings $Y_1$ back to $0.085\%$, at the
cost of $Y_2$ and $Y_3$. Coordinate search can adjust each component separately, whereas the
gradient construction uses one exponent for all four.
